\cotefile{71}
\scanpage{36}\mspage{$\underline{\overline{\mathrm{IX}}}$}
Comme $\Gamma$ est un $J$-polygone et $\partial E(\Gamma)$ est un polygone convexe il est possible de trouver un disque $D_s$ qui vérifie b) de la def.~1 par rapport à $\Gamma$ et à $\partial E(\Gamma)$. Par le lemme $\underline{\overline{\mathrm{IX}}}$ on sait que l'angle $\alpha_1$ du secteur angulaire interieur (de $D_s$) determiné par $\partial E(\Gamma)$ est inferieur à $\pi$, \deleted{alors nécessairement $\alpha_1<\pi$}, car $s$, $s'$ et $s''$ ne sont pas colineaires. Par conséquent l'angle $\alpha_2$ qui sous-tend le secteur angulaire exterieur vérifie $\alpha_2>\pi$.

Mais $\operatorname{ext}\partial E(\Gamma)\subset\operatorname{ext}\Gamma$ implique qu'on peut trouver dans $D_s$ un angle $\alpha'_2>\pi$ qui sous-tend un secteur angulaire contenu dans $\operatorname{ext}\Gamma$, par conséquent le secteur angulaire interieur de $D_s$ \deleted{par} donné par $\Gamma$ a un angle $<\pi$.

Mais \unsure{tout} polygone convexe a au moins trois sommets et cela suffit à prouver le lemme.

On est en mésure de montrer le :

\uline{Lemme X}I : Si $\Gamma$ est un $J$-polygone \inserted{à $n\unsure{>|\geq}3$ \unsure{côtés|cotés}} il existe un couple de sommets $s,s'\in S(\Gamma)$ tels que $s'$ est visible de $s$ de l'interieur de $\Gamma$.
% slanted
\margin{faux pour $n=3$!}

\margin{OPS $\Gamma$ non convexe (sinon c'est clair \unsure{grâce à --})}

\uline{Démonstration} : Soit $s$ un sommet satisfaisant la propriété ennoncée au lemme $\underline{\overline{\mathrm{\unsure{VIII|VII}}}}$%
% slanted
\margin{\uline{lemme X}}, soient $s'$, $s''$ les deux sommets \deleted{consecutifs} \inserted{\unsure{adjacents}} à $s$ et $T$ le tri\uline{an}gle%
% slanted
\margin{ambigu} de sommets $s,s',s''$. Si $(T-([s,s']\cup[s,s'']))\cap\Gamma=\unsure{\varnothing|\phi}$ on a que $]s',s''[\cap\Gamma=\varnothing$, donc $]s',s''[\subset\operatorname{int}\Gamma$ ou bien $]s',s''[\subset\operatorname{ext}\Gamma$. Mais $\inserted{]s',s''[\subset}T-([s,s']\cup[s,s''])\subset\deleted{\unsure{\operatorname{ext}}}\operatorname{int}\Gamma$ (puisqu'il est connexe et et son intersection avec $\operatorname{int}\Gamma$ est non vide par le choix de $s$) donc $]s',s''[\subset\operatorname{int}\Gamma$.
% slanted
\margin{\uline{faux}}
% Pencil leader from the underlined faux.
\margin{il faut prendre $s$ comme dans \unsure{la|le} \uline{\unsure{dém|dém.}} du lemme $\underline{\overline{\mathrm{VIII}}}$ (\unsure{pas son énoncé}!)}

% slanted
\margin{\deleted{\illegible{2}}!}
\uline{Si} \uline{$(T-([s,s']\cup[s,s'']))\cap\Gamma\mathrel{\unsure{\big/\;\subset|\not\subset}}]s',s''[$} on doit pouvoir trouver au moins un sommet $z\in]s',s''[$ (puisque $\Gamma\cap]s',s''[\ne]s',s''[$). On a aussi $\mathring{T}\cap\Gamma=\varnothing$ donc $]s,z[\subset\text{\tikz[baseline=(ic.base)]{\node[inner xsep=3pt,inner ysep=2pt] (ic) {$\mathring{T}\subset\operatorname{int}\Gamma$};\draw (ic.center) ellipse[x radius=1.02cm,y radius=.40cm];\draw ([yshift=-.40cm]ic.center) to[bend right=23] ++(.25,-.24);}}$ (puisque $\mathring{T}\subset\mathbb{R}^2-\Gamma$ et $\mathring{T}\cap\operatorname{int}\Gamma\ne\varnothing$)
% slanted
\margin{\tikz[baseline=(m.base)]{\node[inner sep=1pt] (m) {\shortstack{\unsure{mais|est} non\\vide?}};\draw (m.center) ellipse[x radius=.87cm,y radius=.47cm];}}
% Pencil oval around the second inclusion and a leader from the oval to the note.
% slanted
\margin{dire avant! (et \unsure{tu auras lemme|tu as un lemme|tu as au lemme} $\underline{\overline{\mathrm{VIII}}}$)}
% archive sheet number: 35
\scanpage{37}\mspage{$\underline{\overline{\mathrm{X}}}$}
\margin{!!}Reste le cas $\mathring{T}\cap\Gamma\ne\unsure{\varnothing|\phi}$. Considerons l'enveloppe convexe\margin{ce n'est \unsure{pas une} ``\unsure{env. conv.}''!} $\unsure{\mathbb{C}|C}$ de $(\mathring{T}\cup\inserted{\bigl[}]s',s''[\inserted{\bigr]})\cap\Gamma$.

On a $\mathbb{C}\subset T-([s,s']\cup[s,s''])$. Observons que $S(\partial C)\deleted{\unsure{\subset|\cap}}\cap\mathring{T}\subset S(\Gamma)$.

Comme $\mathring{T}\cap\Gamma\ne\varnothing$ il y a\deleted{\illegible{1}} au moins un sommet de $\partial C$ à l'interieur $\mathring{T}$. \unsure{Le|le} lemme $\underline{\overline{\mathrm{IV}}}$ montre qu'il existe $y\unsure{\notin|\in}S(\partial C)$, $y\in\mathbb{C}$, tel que $y$ est $(\mathbb{R}^2-\mathbb{C})$-visible de $s$

(Noter aussi que $\mathring{C}\ne\varnothing$ puisque $\Gamma\cap\mathring{T}\ne\varnothing\Rightarrow\operatorname{card}(\Gamma\cap]s',s''[)\geq 2$)

Soit $\alpha\in A(\partial C)$ telle que $y\in\alpha$ le lemme V montre que tout point de $\alpha$ est $(\mathbb{R}^2-C)$-visible de $s$. Si $\alpha=[z,z']$ montrons que $z^{\deleted{\unsure{\prime}}}\notin[s',s'']$ ou bien $z^{\unsure{\prime|\prime\prime}}\notin[s',s'']$. Soit donc $z''\in\mathring{T}\cap\Gamma$ et considerons le triangle $T'$ de sommets $z,z',z''$. Si $z',z''\in[s',s'']$ on voit que $]\unsure{s|S},y[$ ($y\in]z,z'[$) coupe $\mathring{T}'$ selon un ensemble non-vide. Cela contredit le choix de $y$ puisque $\mathring{T}'\subset\mathbb{C}$, \deleted{\unsure{C.a|c.a|c.à}} donc $z'$ par exemple vérifie $z'\notin[s',s'']$, $z'\in\unsure{S'|S}(\Gamma)$ (puisque $z'\in S(\partial\mathbb{C})\cap\mathring{T}$) et le lemme \deleted{\unsure{3|5}} V montre que $z'$ est $(\mathbb{R}^2-\mathbb{C})$-visible de $s$, donc $\operatorname{int}\Gamma$-visible de $s$, c.q-f.d.
% margin note written sideways
\margin{compliqué, prendre \unsure{un|le|les} sommet de $\Gamma$ \unsure{$\in$} \deleted{\unsure{$\mathring T$|$\alpha$}} \unsure{et} qui est à dist. max. de la droite \deleted{\unsure{$\Delta$}} $(s's'')$}

\uline{Corollaire à la demonstration du lemme $\underline{\overline{\mathrm{\unsure{VIII|VII}}}}$.} \margin{!!}

Si $\Gamma$ est un polygone il existe un couple de sommets \inserted{\unsure{distincts}} $s,s'\in S(\Gamma)$ tels que $s'$ est \deleted{\unsure{vis}} $(\mathbb{R}^2-\Gamma)$-visible de $s$ et $]s',s[\subset\partial E(\Gamma)$.
% archive sheet number: 36
\scanpage{38}\mspage{$\underline{\overline{\mathrm{XI}}}$}
3 - \uline{Preuve du \unsure{théorĕme|théorème}} 1

\margin{!}Soit $\Gamma$ un polygone et $E(\Gamma)$ son enveloppe convexe. Soit $s,s'\in S(\Gamma)$ un couple de sommets \deleted{\illegible{1}} tels que $s'$ est $(\mathbb{R}^2-\Gamma)$-visible de $s$ et $]s,s'[\subset\partial E(\Gamma)$ (cf. coroll. au lemme $\underline{\overline{\mathrm{VIII}}}$). Soient $\Gamma^1_{ss'}$ et $\Gamma^2_{ss'}$ les deux arcs de polygone \inserted{(de $\Gamma$)} réliant $s$ à $s'$ (cf. Remarque 1).
\margin{$\vert$}

Soit $\Lambda_{ss'}$ l'arc de polygone $[s,s']$. On appelle $\Lambda^1$ (resp. $\Lambda^2$) le polygone obtenu en récollant $\Lambda_{ss'}$ avec $\Gamma^1_{ss'}$ (resp. $\Gamma^2_{ss'}$) selon $\{s,s'\}$.

% slanted
\margin{\shortstack{$a=s$\\$b=s'$}\,?}
\uline{Lemme} $\underline{\overline{\mathrm{XII}}}$. Si $\Lambda^1$ (resp. $\Lambda^2$) est un $J$-polygone alors $\Gamma^2_{ab}-\{a,b\}\subset\operatorname{int}\Lambda^1$ ou bien $\Gamma^2_{ab}-\{a,b\}\subset\operatorname{ext}\Lambda^{\deleted{\unsure{2}}1}$ (resp. $\Gamma^1_{ab}-\{a,b\}\subset\operatorname{int}\Lambda^2$ $\Gamma^1_{ab}-\{a,b\}\subset\operatorname{ext}\Lambda^2$)
\margin{1}

\uline{Demonstration} : Il est clair que $\Gamma^2_{ab}-\{a,b\}\subset\deleted{\unsure{\partial|E}}\mathbb{R}^2-\Lambda^1$ car $]s,s'[\cap\Gamma=\unsure{\varnothing|\phi}$, mais en plus $\Gamma^2_{ab}-\{a,b\}$ est connexe donc la conclusion. On réfait l'argument pour $\Lambda^2$, c.q.f.d.
\margin{$\vert$}\margin{$\vert$}

On suppose maintenant \inserted{que} $\Lambda^1$ et $\Lambda^2$ \deleted{\unsure{que}} sont des $J$-polygones.

\uline{Lemme} $\underline{\overline{\mathrm{XIII}}}$ : Si $x\in\partial E(\Gamma)$ et $D_x$ est un disque de centre $x$ \unsure{vérifiont|vérifiant} b) de la definition 2 par rapport à $\partial E(\Gamma)$ on a alors :
\[
\operatorname{ext}\partial E(\Gamma)\cap\text{\tikz[baseline=(dx.base)]{\node[inner sep=2pt] (dx) {$\deleted{D_x}$};\draw ([xshift=-2pt,yshift=2pt]dx.north west) -- ([xshift=2pt,yshift=-4pt]dx.south east);\draw ([xshift=2pt,yshift=4pt]dx.north east) -- ([xshift=-1pt,yshift=-4pt]dx.south west);}}\subset\operatorname{ext}\Lambda^i,\qquad i=1,2.
\]
\margin{\resizebox{\marginparwidth}{!}{\tikz[baseline=(ex.base)]{\node[inner sep=3pt] (ex) {$\begin{aligned}
\operatorname{Ext}(\partial E(\Gamma))&=\mathbb{R}^{\unsure{2|7}}-E(\Gamma)\\
&\subset\operatorname{Ext}(\Lambda^i)
\end{aligned}$};
\draw (ex.north east) to[bend right=3] (ex.north west) -- (ex.south west) to[bend left=3] (ex.south);}}}

\uline{Démonstration} : Comme $E(\Lambda^i)\subset E(\Gamma)$, $i=1,2$, on a $\operatorname{ext}\deleted{\partial E(}\Lambda^i\deleted{)}\supset\operatorname{ext}\partial E(\Lambda^i)\supset\operatorname{ext}\partial E(\Lambda)$ ce qui entraîne l'ennoncé.
% archive sheet number: 37
\scanpage{39}\mspage{\deleted{\unsure{X bis|XI}}$\underline{\overline{\mathrm{\unsure{XII}}}}$}
On observe que $\Lambda^1$ et $\Lambda^2$ ont une arète en commun\inserted{, savoir,} $[s,s']$ et $\Lambda^1\cap\Lambda^2=[s,s']$.

\uline{Lemme} XIV. Si $D_s$ est un disque satisfaisant b) de la definition 2 par rapport à $\Lambda^1$ et à $\Lambda^2$ on a que $\operatorname{int}\Lambda^1\supset\Lambda^2\cap D_s-[s,s']\cap D_s$ ou bien $\operatorname{int}\Lambda^2\supset\Lambda^1\cap D_s\unsure{-}[s,s']\cap D_s$ \deleted{\unsure{(}}et resp. \inserted{que} $\operatorname{ext}\Lambda^2\supset\Lambda^1\cap D_s-[s,s']\cap D_s$ ou bien $\operatorname{ext}\Lambda^1\supset\Lambda^2\cap D_s-[s,s']\cap D_s$\deleted{)}.

\uline{Demonstration} : D'aprés le lemme XIII, $\operatorname{ext}\Lambda^1\cap D_s\supset\operatorname{ext}\partial E(\Gamma)\cap D_s$, $\operatorname{ext}\Lambda^2\cap D_s\supset\operatorname{ext}\partial E(\Gamma)\cap D_s$ donc $\operatorname{ext}\Lambda^1\cap\operatorname{ext}\Lambda^2\cap D_s\ne\unsure{\varnothing|\phi}$.

Comme $\Lambda^1\cap D_s$ et $\Lambda^2\cap D_s$ ont un coté en \unsure{commun|comun} il résulte que le \inserted{secteur angulaire} $\operatorname{ext}\Lambda^1\cap D_s$ contient ou est contenu dans l'un des secteurs angulaires determinés par $\Lambda^2\cap D_s$. Mais $\operatorname{ext}\Lambda^1\cap\operatorname{ext}\Lambda^2\cap D_s\ne\varnothing$ entraîne que ce secteur est précisement $\operatorname{ext}\Lambda^2\cap D_s$, cela implique (sans peine)
\[
\operatorname{int}\Lambda^1\supset\Lambda^2\cap D_s-[s,s']\cap D_s
\quad\text{ou bien}\quad
\deleted{\operatorname{ext}}\operatorname{int}\Lambda^2\supset\Lambda^1\cap D_s-[s,s']\cap D_s,
\]
c.q.f.d.

La conjonction du lemme $\underline{\overline{\mathrm{XII}}}$ et du lemme $\underline{\overline{\mathrm{XIII}}}$ \deleted{\unsure{entraînent|entraîne}} montre que
\[
\begin{aligned}
\operatorname{int}\Lambda^1\supset{}&\Lambda^2-\{\unsure{g|s},s'\}
&&\text{ou bien }\operatorname{int}\Lambda^2\supset\Lambda^1-\{s,s'\}\ \text{et on a}\\[-2pt]
&\inserted{\substack{\text{''}\\\Gamma^2_{s,s'}-\{s,s'\}}}
&&\inserted{\unsure{-}\qquad\substack{\text{''}\\\Gamma^{\unsure{1|2}}_{s,s'}\unsure{.}-\{s,s'\}}}\\
\text{respectivement}\quad&\operatorname{ext}\Lambda^2\supset\Lambda^1-\{s,s'\}
&&\text{et }\operatorname{ext}\Lambda^1\supset\Lambda^2-\{s,s'\},
\end{aligned}
\]

Supposons pour fixer les idées que $\operatorname{int}\Lambda^1\supset\Lambda^2-\{s,s'\}$ donc aussi que $\operatorname{ext}\Lambda^2\supset\Lambda^1-\{s,s'\}$.

On a donc :

\uline{Lemme XV} : $\operatorname{int}\Lambda^2\subset\operatorname{int}\Lambda^1$.
% archive sheet number: 38
\scanpage{40}\mspage{$\underline{\overline{\mathrm{XIII}}}$}
\uline{Demonstration} : Pour voir que $\operatorname{int}\Lambda^2\subset\deleted{\unsure{\operatorname{ext}}}\operatorname{int}\Lambda^1$ \inserted{on} observe d'abord que :
\[
\Lambda^1-\{s,s'\}\subset\operatorname{ext}\Lambda^2\quad\text{donc}\quad\deleted{\unsure{\operatorname{ext}}}\operatorname{int}\Lambda^2\cap(\Lambda^1-\{s,s'\})=\unsure{\varnothing|\phi}
\]
, et par conséquent :
\[
\operatorname{int}\Lambda^2\subset\operatorname{int}\Lambda^1\quad\text{ou bien}\quad\operatorname{int}\Lambda^2\subset\operatorname{ext}\Lambda^1\quad(\text{car }\operatorname{int}\Lambda^2\subset\mathbb R^2-\Lambda^1).
\]
Mais $\operatorname{int.}\Lambda^2\cap\operatorname{int}\Lambda^1\ne\varnothing$. En effet soit $x\in{]s,s'[}$ et $D_x$ un disque de centre $x$ qui vérifie b) de la def.1 par rapport à $\Lambda^1$, $\Lambda^2$ et $\partial E(\Gamma)$, le lemme $\underline{\overline{\mathrm{XIII}}}$ montre que $D_x\cap\operatorname{ext}\partial E(\Gamma)=D_x\cap\operatorname{ext}\Lambda^1=D_x\cap\operatorname{ext}\Lambda^1$ donc $D_x\cap\deleted{\unsure{\operatorname{ext}}}\operatorname{int}\Lambda^1=D_x\cap\deleted{\unsure{\operatorname{ext}}}\operatorname{int}\Lambda^2$, et par conséquent $\operatorname{int}\Lambda^1\cap\operatorname{int}\Lambda^2\ne\varnothing$, c.q.f.d.

On veut prouver que $\Gamma$ est un $J$-polygone par récurrence sur $\operatorname{card}S(\Gamma)-2$.

Si $\operatorname{card}S(\Gamma)-2=1\quad\Rightarrow\quad\Gamma={}$ au bord d'un triangle\inserted{et le theorème 1 devient donc trivial}, donc il suffit de prouver que si le \unsure{théorĕme|théorème} 1 est vrai pour tout $\Gamma$ avec $\operatorname{card}S(\Gamma)-2<n$ alors il est aussi vrai pour tout $\Gamma$ avec $\operatorname{card}S(\Gamma)-2=n$.

Il est clair que $\operatorname{card}S(\Lambda^i)<\operatorname{card}S(\Gamma)$ \inserted{$i=1,2$}, autrement $\Lambda^i=\Gamma$ ce qui est absurd puisque $ {]s,s'[}\cap\Gamma=\varnothing$. Les lemmes $\underline{\overline{\mathrm{XII}}}$ et $\underline{\overline{\mathrm{XIV}}}$ \inserted{avec le lemme $\underline{\overline{\mathrm{XV}}}$} nous permettent de supposer sans perte de generalité que : a) $\operatorname{int}\Lambda^2\subset\operatorname{int}\Lambda^1$ ;

b) $\Lambda^2-\{s,s'\}=\Gamma^2_{s,s'}-\{s,s'\}\subset\operatorname{int}\Lambda^1$.

Posons donc :
\[
\begin{aligned}
V&=\operatorname{ext}\Lambda^1\cup\operatorname{int}\Lambda^2\cup{]s,s'[}\quad\text{et}\\
U&=\operatorname{int}\Lambda^1-(\operatorname{int}\Lambda^2\cup(\Gamma^2_{ss'}\cap\operatorname{int}\Lambda^2)).
\end{aligned}
\]
Par \unsure{hypothèse|hypothēse} inductive $\operatorname{ext}\Lambda^1$ est non-relativement-compact donc $V$ non plus et $U\subset\operatorname{int}\Lambda^1$ est relativement compact puisque $\operatorname{int}\Lambda^1$ est relativement compact par hypothèse de récurrence.

\figure{\makebox[\linewidth][r]{\begin{tikzpicture}[x=1cm,y=1cm]
\fill (0,0) -- (.11,.015) -- (.15,-.06) -- cycle;
\draw (.05,-.015) .. controls (.94,-.34) and (1.60,-.80) .. (3.8,-1.20);
\end{tikzpicture}}}
% archive sheet number: 39
\scanpage{41}
% Two lines written upside down relative to the archive sheet number.
Demonstration : pour voir que $\operatorname{int}\Lambda^2\supset\operatorname{int}\Lambda^1$ il suffit de prouver que tout point \unsure{de|d}
% archive sheet number: 40
\scanpage{42}\mspage{$\underline{\overline{\mathrm{XIV}}}$}
\uline{Lemme} $\underline{\overline{\mathrm{XVI}}}$. On a : a) $\mathbb R^2-\Gamma=U\cup V$

b) $U$ et $V$ sont des ouverts connexes, $U$ relativement compact, $V$ non-relativement compact.

\uline{Démonstration} : a) On a $U\cup V=\operatorname{ext}\Lambda^1\cup\operatorname{int}\Lambda^1\cup{]s,s'[}-(\Gamma^2_{ss'}-\{s,s'\})$.

D'autre côté $\Gamma=\Lambda^1-{]s,s'[}\cup(\Gamma^2_{ss'}-\{s,s'\})$, mais alors
\[
\begin{aligned}
\mathbb R^2-\Gamma&=(\mathbb R^2-\Lambda^1)\cup{]s,s'[}-(\Gamma^2_{ss'}-\{s,s'\})=\operatorname{ext}\Lambda^1\cup\operatorname{int}\Lambda^1\cup{]s,s'[}\\
&\hspace{12em}-(\Gamma^2_{ss'}-\{s,s'\}).
\end{aligned}
\]

b) Prouvons que $U$ est connexe. Soient $x,y\in U$\unsure{;|,} montrons que $x$ peut être relié\deleted{\unsure{r|e}} à $y$ par un arc de polygone \unsure{.}contenu dans $U$. Si $x,y\in U$ alors $x,y\in\operatorname{int}\Lambda^1$.

On peut donc reli$\overset{\unsure{\prime|\cdot}}{\text{e}}$r $x$ à $y$ par un arc de polygone $\Lambda_{xy}\subset\operatorname{int}\Lambda^1$. Si $\Lambda_{xy}\subset\underbrace{\operatorname{int}\Lambda^1-R_{\Lambda^2}}_{\substack{\text{''}\\\underset{\unsure{\cdot}}{U}}}$ $\bigl(R_{\Lambda^2}=\operatorname{int}\Lambda^2\cup\Lambda^2\bigr)$ il n'y a plus rien à montrer.

Si $\Lambda_{xy}\cap(\Gamma^2_{ss'}-\{s,s'\})=\unsure{\varnothing|\phi}$ alors $\Lambda_{xy}\cap\operatorname{int}\Lambda^2=\varnothing$, puisque $x,y\in\operatorname{ext}\Lambda^2$ et $\Lambda_{xy}\cap\Lambda^2=\varnothing$ (car $\Lambda^2=[s,s']\cup(\Gamma^2_{ss'}-\{s,s'\})$ et $\Lambda_{xy}\cap[s,s']=\varnothing$ puisque $\Lambda_{xy}\subset\operatorname{int}\Lambda^{\unsure{1|4}}$)

Si $\Lambda_{\unsure{x}y}\cap(\Gamma^2_{ss'}-\{s,s'\})\ne\varnothing$ on peut trouver $x'$ et $y'\in(\Gamma^2_{ss'}-\{s,s'\})\cap\Lambda_{xy}$ tels que $\Lambda_{xx'}-\{x'\}\subset\operatorname{ext}\Lambda^2$, $\Lambda_{y'y}-\{y'\}\subset\operatorname{ext}\Lambda^2$ (ici $\Lambda_{xx'}$ (resp. $\Lambda_{y'y}$) dénote l'arc de polygone obtenu à partir de $\Lambda_{xy}$ qui rélie $x$ a $x'$ (resp. $y'$ à $y$) \inserted{obtenu} de \unsure{foçon|façon} analogue a \unsure{$\overset{\unsure{\mathrm c|2|?}}{\text{ce}}$lui} de la Remarque 1), pour cela il suffit d'avoir
\[
\tikz[baseline=(eq.base)]{\node[inner sep=3pt] (eq) {$\begin{aligned}
\Lambda_{xx'}\cap(\Gamma^2_{ss'}-\{s,s'\})&=x',\\
\Lambda_{y y'}\cap(\Gamma^2_{ss'}-\{s,s'\})&=y'.
\end{aligned}$};
\draw ([xshift=2pt]eq.north west) -- (eq.south west) -- (eq.south east) -- ([xshift=4pt]eq.north east);
\node[inner sep=.2pt] at ([xshift=9pt]eq.east) {$\unsure{*}$};}
\]

\deleted{\unsure{Soit maintenant $\Gamma$}} \inserted{Comme} On a $x',y'\in\Gamma^2_{ss'}-\{s,s'\}$ on peut trouver un arc de polygone $\Gamma_{x'y'}$ de $\Gamma$, $\Gamma_{x'y'}\subset\Gamma^2_{ss'}-\{s,s'\}$ (cf. Remarque 1).
% archive sheet number: 41
